% ---------------------------------------------------------------------------
% Estilos compartilhados pelas figuras adaptadas de Andrade et al. (2019).
%
% As duas `pics` são reproduzidas literalmente de styles.tikz, do material
% original dos autores. ATENÇÃO: em `Machine`, o rótulo é desenhado por
% `node[midway] {#10}` -- isto é, o argumento #1 seguido de um `0` literal.
% Logo `lenght=10` imprime "100" e `lenght=3` imprime "30". O argumento
% `label=` declarado nunca é usado. Reescrever isso mudaria silenciosamente
% todos os números de todas as figuras; alteramos apenas as cores.
% ---------------------------------------------------------------------------

% Vermelho do "X" de inviabilidade (ATTFunctionalRed no original).
\definecolor{myred}{RGB}{207, 42, 42}

% Máquinas (figura da instância): as cores identificam a máquina.
\colorlet{MaqUm}{mylightorange1}
\colorlet{MaqDois}{mylightblue1}

% Tarefas (figuras do decodificador e do escalonamento): as cores
% identificam a tarefa.
\colorlet{JobALight}{mylightorange1}
\colorlet{JobADark}{mydarkorange1}
\colorlet{JobBLight}{mylightblue1}
\colorlet{JobBDark}{mydarkblue1}

% Revelação por overlay mantendo a caixa delimitadora estável: os elementos
% ocultos continuam ocupando espaço, então a figura não "salta" entre os
% passos. Mesma técnica usada na apresentação do MIC 2026.
\tikzset{%
  invisible/.style={opacity=0},
  visible on/.style={alt={#1{}{invisible}}},
  alt/.code args={<#1>#2#3}{%
    \alt<#1>{\pgfkeysalso{#2}}{\pgfkeysalso{#3}}%
  },
}

\tikzset{%
  pics/Machine/.style args={lenght=#1, color=#2, label=#3}{code={%
      \draw[inner sep=0mm, minimum size=5mm, rounded corners=1.5mm,fill=#2]
         (0.1, 0.1) rectangle ++($ (#1 - 0.2, 0.8) $);
    \path (0,0) --  (#1,1) node[midway] {#10};
  }},%
  %
  pics/MachineCap/.style args={lenght=#1}{code={%
      \draw[inner sep=0mm, minimum size=5mm, rounded corners=1mm,
            fill=mygray]
         (0.0, 0.0) rectangle ++($ (#1 + 0.1, 1.1) $);
  }}%
}
